\documentclass[12pt]{article}
\usepackage[T1]{fontenc}
\usepackage[utf8]{inputenc}
\usepackage{lmodern}
\usepackage{textcomp}
\usepackage{enumitem}
\usepackage{geometry}
\geometry{margin=1in}
\begin{document}
\begin{center}
Answer Key - Economics - Undergraduate Level Assessment 23 \
\small (Answers are ordered exactly as in the question paper.)
\end{center}

\section*{Multiple Choice}
\begin{enumerate}[label=\arabic*.]
\item \textbf{Answer:} A
\\ \textit{Options:} A) Economic resources are allocated according to the decisions of the central bank.; B) Private property is fundamental to innovation, growth, and trade.; C) A central government plans the production and distribution of goods.; D) Most wages and prices are legally controlled.
\\ \textit{Note:} The correct answer reflects the influence of a central bank in a capitalist market economy, where monetary policy decisions help regulate economic activity and resource allocation through interest rates and money supply, rather than through direct government intervention.
\item \textbf{Answer:} D
\\ \textit{Options:} A) Seasonal unemployment; B) Casual unemployment; C) Cyclical unemployment; D) Structural unemployment; E) Underemployment
\\ \textit{Note:} This is an example of structural unemployment because the technician's job loss is due to a fundamental change in the market demand for film photography, leading to a mismatch between his skills and the current needs of the economy.
\item \textbf{Answer:} D
\\ \textit{Options:} A) The largest and smallest; B) The smallest and second smallest; C) The middle two; D) The smallest 2; E) The largest 2
\\ \textit{Note:} In the Johansen "trace" test, the smallest eigenvalues correspond to the number of cointegrating vectors under the null hypothesis, and since the test is assessing for 2 cointegrating vectors in a 4-variable system, the smallest 2 eigenvalues are used to determine if the null hypothesis can be rejected.
\item \textbf{Answer:} A
\\ \textit{Options:} A) 75 percent.; B) 0.5 percent.; C) 20 percent.; D) 100 percent.; E) 25 percent.
\\ \textit{Note:} The correct answer is 75 percent because economic growth is calculated using the formula: ((GDP per capita in 2000 - GDP per capita in 1990) / GDP per capita in 1990) x 100, which in this case is ((15000 - 10000) / 10000) x 100 = 50\%.
\item \textbf{Answer:} C
\\ \textit{Options:} A) An increase in consumer confidence.; B) Increase in the prices of goods and services.; C) Business firms expect lower sales in the future.; D) A technological advancement in the production process.; E) An increase in the wealth of consumers.
\\ \textit{Note:} Business firms expecting lower sales in the future indicates a decrease in business investment and consumer spending, which collectively reduce aggregate demand, leading to a leftward shift of the aggregate demand curve.
\end{enumerate}

\section*{True / False}
\begin{enumerate}[label=\arabic*.]
\setcounter{enumi}{5}
\item \textbf{Answer:} True
\\ \textit{Options:} A) True; B) False
\\ \textit{Note:} Theo's job loss due to seasonal pool closure is classified as natural unemployment because it results from predictable and unavoidable seasonal fluctuations in demand for labor, rather than from economic downturns or structural changes in the labor market.
\item \textbf{Answer:} False
\\ \textit{Options:} A) True; B) False
\\ \textit{Note:} The statement is false because higher wages generally incentivize workers to increase their productivity rather than decrease it, as they feel more motivated and valued in a competitive labor market.
\item \textbf{Answer:} True
\\ \textit{Options:} A) True; B) False
\\ \textit{Note:} The velocity of money is calculated by dividing the total transaction value by the money supply, so with a transaction value of $100,000 and a money supply of $10,000, the velocity is 100,000 / 10,000 = 10, not 20, making the statement false.
\item \textbf{Answer:} False
\\ \textit{Options:} A) True; B) False
\\ \textit{Note:} The statement is false because coefficient estimates in Vector Autoregressions (VARs) can be both positive and negative, indicating that the relationships among the variables can be either direct or inverse depending on the dynamics of the system being analyzed.
\item \textbf{Answer:} True
\\ \textit{Options:} A) True; B) False
\\ \textit{Note:} The statement is true because when a good has many close substitutes, consumers can easily switch to those alternatives in response to price changes, leading to a more significant change in the quantity demanded.
\end{enumerate}

\section*{Long Form Response}
\begin{enumerate}[label=\arabic*.]
\setcounter{enumi}{10}
\item \textbf{Answer:} Multicollinearity significantly impacts the properties of the Ordinary Least Squares (OLS) estimator by compromising its efficiency and potentially introducing bias in the estimation process. When independent variables are highly correlated, it becomes difficult to isolate the individual effect of each variable on the dependent variable, leading to inflated standard errors. This inflation reduces the reliability of the coefficient estimates, making them less efficient in terms of variance. Although the OLS estimator remains unbiased in the presence of multicollinearity, the lack of precision in estimating coefficients can lead to misleading conclusions about the relationships being studied. Therefore, while OLS estimators may provide unbiased estimates, their inefficiency due to multicollinearity undermines the robustness of the regression results.
\\ \textit{Note:} correct because it accurately describes how multicollinearity inflates standard errors and diminishes the efficiency of the OLS estimator, ultimately affecting the reliability of coefficient estimates by making it challenging to determine the individual effects of correlated independent variables.
\item \textbf{Answer:} A kinked demand curve characterizes oligopoly markets where firms face a demand curve that is more elastic for price increases and less elastic for price decreases. This results in a kink at the current market price, leading to price rigidity; firms are reluctant to raise prices due to the anticipated loss of customers and are also hesitant to lower prices since competitors will match this decrease. The condition P \textless{} MC does not align with the principles of the kinked demand curve because, in equilibrium, firms typically set prices above marginal costs to maximize profits. If a firm were to set P \textless{} MC, it would incur losses, contradicting the purpose of profit maximization in oligopoly behavior. Thus, the kinked demand curve illustrates how firms maintain stable pricing despite fluctuating costs and demand conditions.
\\ \textit{Note:} correct because it accurately describes how the kinked demand curve leads to price rigidity in oligopoly markets, where firms' pricing behavior is influenced by the differing elasticities of demand for price increases and decreases, and it clarifies that the condition P \textless{} MC contradicts the equilibrium established by firms' strategic pricing decisions.
\end{enumerate}

\vfill
\noindent\textit{End of Answer Key}
\end{document}